\section{State-by-State Comparison: Tax Planning for Different Project Scenarios}

\subsection{Texas: Temporary Sales and Use Tax Exemptions}

Under Texas Tax Code \S{}\S{} 151.359 and 151.3595, Texas provides temporary sales and use tax exemptions for qualifying data centers and qualifying large data center projects. A qualifying data center is exempt from the 6.25\% state sales and use tax on qualifying purchases, while local sales and use tax remains due; the exemption period is 10 or 15 years. For a qualifying large data center project, qualifying purchases are exempt from both state and local sales and use tax for 20 years.

The exemption covers electricity and equipment that are necessary and essential to data center operations, including servers, data storage devices and network connectivity equipment. For electricity purchases, a predominant use study is generally required to distinguish taxable from nontaxable electricity use unless the data center or project satisfies the applicable stand-alone facility requirements.

% Transcribed from assets/figure-2.png.
% The identical common requirements are shown once, applying to both types.
\begingroup
\ReportTableSetup
\begin{longtable}{|L{0.28\linewidth}|L{\dimexpr0.72\linewidth-4\tabcolsep-3\arrayrulewidth\relax}|}
\caption{Eligibility Requirements}\label{tab:texas-eligibility}\\
\hline
\multicolumn{2}{|L{\dimexpr\linewidth-2\tabcolsep-2\arrayrulewidth\relax}|}{\textbf{Common Requirements (Both Types)}} \\ \hline
\multicolumn{2}{|L{\dimexpr\linewidth-2\tabcolsep-2\arrayrulewidth\relax}|}{1. Hold or concurrently apply for a Texas Sales or Use Tax Permit or Texas Direct Payment Permit;} \\ \hline
\multicolumn{2}{|L{\dimexpr\linewidth-2\tabcolsep-2\arrayrulewidth\relax}|}{2. Submit documentation outlining the plan to meet capital investment and job-creation requirements, together with a proposed site map or building schematic;} \\ \hline
\multicolumn{2}{|L{\dimexpr\linewidth-2\tabcolsep-2\arrayrulewidth\relax}|}{3. Be located in Texas;} \\ \hline
\multicolumn{2}{|L{\dimexpr\linewidth-2\tabcolsep-2\arrayrulewidth\relax}|}{4. Be specifically constructed or refurbished primarily to house servers and related equipment and support staff for the processing, storage and distribution of data;} \\ \hline
\multicolumn{2}{|L{\dimexpr\linewidth-2\tabcolsep-2\arrayrulewidth\relax}|}{5. Have an uninterruptible power source, generator backup power, a sophisticated fire suppression and prevention system, and enhanced physical security including restricted access, video surveillance and electronic systems;} \\ \hline
\multicolumn{2}{|L{\dimexpr\linewidth-2\tabcolsep-2\arrayrulewidth\relax}|}{6. Not be used primarily by a telecommunications provider to deliver telecommunications services;} \\ \hline
\multicolumn{2}{|L{\dimexpr\linewidth-2\tabcolsep-2\arrayrulewidth\relax}|}{7. Not be subject to an agreement limiting the appraised value of the data center's property under Tax Code Chapter 313.} \\ \hline

\textbf{Type} & \textbf{Specific Requirements} \\ \hline
Qualifying Data Center &
1. Commit to creating at least 20 qualifying jobs in the county;\par
2. Commit to making, or agreeing to make, at least \$200 million in capital investment within five years;\par
3. Have at least 100,000 square feet of building space used by a single qualifying occupant.
\\ \hline
Qualifying Large Data Center Project &
1. Commit to creating at least 40 qualifying jobs in the county;\par
2. Commit to making, or agreeing to make, at least \$500 million in capital investment within five years;\par
3. Have at least 250,000 square feet of building space used by a single qualifying occupant.
\\ \hline
\end{longtable}
\endgroup


A qualifying data center eligible for the exemption must have at least 100,000 square feet of space used by a single qualifying occupant, make at least \$200 million in capital investment within five years after certification, and create at least 20 qualifying jobs. For a qualifying large data center project, the thresholds increase to at least 250,000 square feet of total space, at least \$500 million in capital investment within five years, and at least 40 qualifying jobs.

The implementing requirements expressly allow a data center facility to be specifically constructed or refurbished. Accordingly, the conversion of an existing mining facility is not disqualified merely because of the facility's original intended use.

To continue benefiting from the exemption, the relevant businesses must also comply with ongoing requirements, including meeting the applicable capital investment and qualifying-job commitments, submitting required reports to the Texas Comptroller, and maintaining complete records of exempt purchases until the data center's certification expires. The Comptroller reviews compliance and conducts audits; a data center that fails to satisfy the requirements may lose the exemption or have its certification revoked.

\subsection{Virginia: Retail Sales and Use Tax Exemption, with an Electricity Consumption Tax}

Code of Virginia \S{} 58.1-609.3(18) provides an exemption from the Retail Sales and Use Tax for qualifying computer equipment and enabling software purchased or leased by eligible data center operators and their tenants (the DCRSUT Exemption). The exemption is currently available through June 30, 2035. If a data center satisfies specified capital investment and job-creation requirements, the exemption may be extended to 2040 or 2050 through a Memorandum of Understanding (MOU).

Qualifying computer equipment includes servers, routers and data storage equipment, and expressly covers cabling, switches, wiring and similar items used to support the operation of exempt equipment. However, general building improvements and fixtures unrelated to data processing, such as lighting and fencing, are excluded, as is fuel otherwise subject to retail sales and use tax that is used to provide electricity.

Unlike Texas, data center operators and tenants in general localities need to make at least \$150 million in new capital investment. In addition, the data center must create at least 50 jobs associated with data center operations or maintenance, each paid at least 150\% of the prevailing annual average wage in the locality. In distressed localities, the statutory thresholds are reduced to \$70 million in capital investment and 10 operations- or maintenance-related jobs.

\ReportImage{3}{figure}{Distribution of General Localities (Gray) and Distressed Localities (Blue)}

The DCRSUT Exemption carries recordkeeping and reporting requirements. Data centers and their tenants must maintain complete and accurate records of capital investment, the number of new jobs and their average annual wages, expenditures on qualifying computer equipment or enabling software, and the value of the DCRSUT Exemption. On behalf of itself and its tenants, the data center must submit annual progress reports and a final report on achievement of performance targets to the Virginia Economic Development Partnership (VEDP) and the Virginia Department of Taxation, as well as ongoing annual reports covering capital investment, total data center jobs, average annual wages, equipment expenditures and the value of the exemption.

If a data center fails to meet the minimum capital investment, job-creation or other performance targets within the prescribed period, generally three years, the relevant businesses must pay the full amount of foregone sales and use tax plus interest.

Beginning July 1, 2026, Virginia implemented a new Data Center Electricity Consumption Tax. For a two-year period, data centers located in the Commonwealth are subject to a tax of \$0.011 per kilowatt-hour of electricity consumed.

\subsection{Georgia: State and Local Sales and Use Tax Exemption}

Under Ga. Comp. R. \& Regs. r. 560-12-2-.117, purchases and uses of qualifying High-Technology Data Center Equipment by a High-Technology Data Center and its customers are exempt from state and local sales and use tax.

In addition to the owner of a High-Technology Data Center, the exemption is available to a High-Technology Data Center Customer, meaning a client, tenant, licensee or end user that is party to a qualifying contract with the Data Center Owner.

The exemption covers computer equipment and materials, software, hardware or equipment used to create, manage or maintain the data center environment, including emergency backup generators, power distribution units and routers. Real Property and electricity used in the ordinary course of operations are not included in the exemption.

The main conditions include the following: the equipment for which the exemption is claimed must be incorporated into or used in the relevant High-Technology Data Center; and the purchaser must be a Data Center Owner or High-Technology Data Center Customer holding a certificate of exemption. A Data Center Owner applying for a certificate of exemption is not subject to a minimum floor-area requirement, but must still meet the applicable Minimum Investment Threshold and New Quality Jobs requirements. A High-Technology Data Center Customer's eligibility depends on the validity of the owner's exemption certificate and the contractual relationship between the parties.

% Transcribed from assets/figure-4.png; all values retained.
\begin{table}[htbp]
\ReportTableSetup
\caption{Requirements for Data Center Owners Applying for an Exemption Certificate During the Investment Period (for Applications on or after May 9, 2022)}
\label{tab:georgia-thresholds}
\begin{tabularx}{\linewidth}{|Y|Y|Y|}
\hline
\textbf{County Population} & \textbf{Average Number of New Quality Jobs} & \textbf{Qualifying Aggregate Expenditures (USD)} \\ \hline
More than 50,000 & 25 & \$250 million \\ \hline
30,000--50,000 & 10 & \$75 million \\ \hline
Fewer than 30,000 & 5 & \$25 million \\ \hline
\end{tabularx}
\end{table}


Correspondingly, while benefiting from the exemption, a High-Technology Data Center Owner must also comply with statutory annual reporting, recordkeeping and other requirements.

\subsection{Tax Planning Approaches}

Tax incentives vary across jurisdictions in their eligibility thresholds, scope of exemption, effective period and application procedures. Mining companies therefore need to assess the tax treatment and its impact in light of the specific facts of each project. The following examples illustrate projects with different characteristics:

Mining Company A plans to invest in or retrofit a large data center with more than \$500 million in investment and more than 250,000 square feet of building space. If located in Texas, the project may qualify as a qualifying large data center project, allowing qualifying purchases of equipment and electricity to be exempt from both state and local sales and use tax. The 20-year exemption period and inclusion of electricity purchases within the exemption give Texas a clear advantage for long-term project operations.

Mining Company B plans to convert a mining site in a low-population county, with an investment of less than \$200 million and limited job creation. If located in a low-population county in Georgia, the project may qualify for the tax incentive by meeting thresholds of only \$25 million in investment and five New Quality Jobs. This approach allows the project to secure tax relief based on its own characteristics, but it also faces disadvantages, including the lack of an electricity exemption and a shorter exemption window. If equipment purchases are delayed, the remaining incentive window will shorten accordingly.

Mining Company C plans to retain the campus and power infrastructure while multiple tenants separately purchase servers and networking equipment. If located in Virginia, the state's framework allows tenants to enter the exemption program through a participation arrangement, supporting a more flexible commercial structure. If the project is located in a distressed locality, the \$70 million investment threshold and 10-job requirement may also be attractive. However, the concurrent electricity consumption tax may reduce the benefit of the equipment exemption.

For mining companies in transition, the future duration of the exemption period, whether electricity consumption receives preferential tax treatment, whether tenants are covered by the exemption, whether the project satisfies the applicable eligibility thresholds, and whether the jurisdiction has suspended or is preparing to suspend new exemption applications are all important factors in an overall tax planning analysis.

\FloatBarrier
